\section{Introduction}
\label{sec:introduction}

Three-dimensional human pose estimation from ordinary video has become an
accessible alternative to specialized motion-capture laboratories. Modern
single-image systems can recover a dense articulated body representation under
varied appearance and imaging conditions; SAM 3D Body, for example, predicts a
full-body mesh based on the Momentum Human Rig representation
\citep{Yang2026SAM3DBody}. Nevertheless, a monocular 3D reconstruction remains
view dependent. Depth ambiguity, self-occlusion, motion blur, truncation, and
imperfect scale recovery affect frontal and lateral observations differently.
When two cameras observe the same movement, their predictions therefore contain
both complementary information and incompatible errors.

Most multi-view 3D pose research resolves this problem before or during 3D
estimation. Calibrated approaches combine image features or 2D heatmaps and
triangulate a common 3D pose \citep{Iskakov2019Learnable,Qiu2019CrossView,
Remelli2020Lightweight}. More recent methods relax calibration or 3D-label
requirements, but they still begin from images or 2D observations and jointly
infer cameras, 2D poses, or 3D locations
\citep{Kim2022CrossViewSelfFusion,Usman2022MetaPose,Jiang2023Probabilistic,
Srivastav2024SelfPose3D}. These formulations do not directly address a common
downstream situation: two complete monocular 3D pose streams already exist, the
camera geometry is unavailable or unreliable, and the desired output is a
single temporally coherent 3D sequence.

Naive averaging is inadequate in this setting. The two streams may use
different translations, orientations, and scales, and a method that minimizes
coordinate noise can suppress the rapid trunk rotation that is central to the
movement under study. This distinction matters for gymnastics and other
rotation-dominant actions, where a visually smooth sequence is not necessarily
a faithful sequence. A useful fusion method must therefore reduce view-specific
artifacts while preserving angular range and peak motion.

This paper formulates \emph{post-estimation two-view 3D pose fusion}. The method
receives temporally aligned face-view and side-view 3D keypoints and produces a
full fused sequence without camera parameters, manually labeled 3D poses, or a
triangulated target during training. Each view is canonicalized in a
pelvis-centered, trial-scaled body frame. Pelvis-to-thorax relative rotations,
wrapped axial angle, angular derivatives, per-view reliability, and symmetric
cross-view disagreement are then supplied to a view-swap-invariant temporal
network. The network predicts a bounded correction to a deterministic
quality-weighted mean. Training uses seeded synthetic corruptions together with
consensus, skeletal, temporal, and complete-cycle rotation constraints.

The study is motivated by the following research question: \emph{Can two
temporally aligned monocular 3D pose streams be fused without camera calibration
or 3D supervision so that cross-view artifacts decrease without attenuating the
trunk-rotation dynamics of a complete movement cycle?} We evaluate two distinct
properties rather than collapsing them into one claim. First, agreement with a
triangulated stream is reported as agreement with a shared upstream
pseudo-reference. Second, reference-independent diagnostics quantify corruption
recovery, skeletal consistency, view-swap invariance, temporal smoothness, and
rotation preservation. Because no independent motion-capture measurements are
available, the study does not claim absolute positional or biomechanical
accuracy. A limited synthetic benchmark against Unity native 3D separately
tests transfer outside the private videos and exposes the boundary between
uncalibrated direct-3D fusion and calibrated 2D triangulation.

A secondary application question asks whether person-disjoint A6 outputs can
support two complementary comparisons: differences between the
elderly-labeled and student-labeled participant groups, and differences among
repeated cycles within the same person. This analysis is observational. The
cohort labels do not identify causal ageing effects, and the fused descriptors
are not clinical measurements.

The contributions are as follows:
\begin{itemize}
  \item We define a practical post-estimation fusion problem in which two
  monocular 3D pose streams, rather than images or 2D heatmaps, are the inputs.
  \item We introduce a rotation-aware, view-swap-invariant residual temporal
  model that combines canonical 3D coordinates with explicit pelvis--thorax
  rotation and deterministic quality features.
  \item We construct a self-supervised objective from reproducible corruptions,
  reliable cross-view consensus, skeletal structure, adaptive temporal dynamics,
  and complete-cycle axial range, while isolating triangulated poses from model
  training and checkpoint selection.
  \item We provide an auditable evaluation protocol that keeps the deterministic
  results, the 14-person held-out learned comparison, descriptive all-person
  diagnostics, and Unity native-3D evidence separate.
  \item We demonstrate a downstream person-disjoint analysis that separates
  adjusted cohort contrasts, within-person cycle variability, and repetition
  trends, with pose-source and limited seed sensitivity checks.
\end{itemize}

The primary learned comparison uses the frozen 14-person test split; the
96-person training and 27-person validation sets are not reused for its
generalization claim. Inference over all 137 people is descriptive only.
Downstream A6 analysis uses ten person-disjoint outer folds and is qualified by
the same centered mixed-model estimand applied to single-view and deterministic
pose sources. These boundaries keep pseudo-reference agreement, limited
synthetic external validity, motion preservation, and observational group
comparisons distinct.
